\section*{บทคัดย่อ}
บทความนี้นำเสนอการออกแบบและพัฒนาระบบลานิสิตสำหรับคณะวิศวกรรมศาสตร์ กำแพงแสน มหาวิทยาลัยเกษตรศาสตร์ เพื่อรวมการยื่นคำขอ การแนบหลักฐาน การพิจารณาตามลำดับ และการติดตามสถานะไว้ในเว็บแอปพลิเคชันเดียว ระบบใช้ Next.js สำหรับส่วนติดต่อผู้ใช้ NestJS สำหรับบริการฝั่งเซิร์ฟเวอร์ และ Prisma เชื่อมต่อฐานข้อมูล การเข้าถึงคำขอและรายงานจำกัดตามตัวตนผู้ใช้ บทบาท และความสัมพันธ์กับนิสิต ระบบข้อมูลสนับสนุน จัดการข้อมูลรายวิชา ผู้สอน และที่ปรึกษา ส่วนคลังข้อมูลรายงานแยกจากข้อมูลธุรกรรมเพื่อรองรับการสรุปจำนวนใบลา วันลา และรายวิชาที่ได้รับผลกระทบ ผลการตรวจซอร์สโค้ดพบองค์ประกอบรองรับกระบวนการหลักและการควบคุมสิทธิ์ดังกล่าว อย่างไรก็ตาม ยังต้องประเมินการใช้งานกับผู้ใช้จริงและทดสอบประสิทธิภาพก่อนสรุปผลเชิงปริมาณ
\par\noindent\textbf{คำสำคัญ:} ระบบลานิสิต, เว็บแอปพลิเคชัน, การอนุมัติ
\section*{Abstract}
This article presents the design and implementation of a student leave management system for the Faculty of Engineering at Kamphaeng Saen, Kasetsart University. The system integrates request submission, evidence attachments, sequential approval, and status tracking in a web application. Next.js provides the user interface, while NestJS and Prisma support server services and database access. Request and report access is scoped by user identity, role, and relationships with students. A separate Student Leave Reference Data Management System manages subject, lecturer, and advisor reference data. The reporting warehouse is separated from operational transactions to summarize requests, leave days, and affected subjects. Source code inspection identifies support for these core workflows and access controls. User evaluation and performance testing remain necessary before quantitative benefits can be established.
\par\noindent\textbf{Keywords:} student leave, web application, approval workflow
